% Part: lambda-calculus
% Chapter: syntax
% Section: terms

\documentclass[../../../include/open-logic-section]{subfiles}

\begin{document}

\olfileid{lam}{syn}{trm}
\olsection{પદો}

લૅમ્બડા કલનનાં પદો ચલોનો અનંત પુરવઠો $\Obj{v_0}$, $\Obj{v_1}$,
\dots{}, સંકેત~``$\lambd$'' અને કૌંસોમાંથી આગમનથી બને છે. ચલો
દર્શાવવા આપણે $x$, $y$, $z$, \dots{} વાપરીશું અને પદો દર્શાવવા
$M$, $N$, $P$, \dots{} વાપરીશું.

\begin{defn}[પદો] \ollabel{defn:term}
લૅમ્બડા કલનનાં \emph{પદો}નો ગણ આગમનથી આ રીતે વ્યાખ્યાયિત થાય છે:
\begin{enumerate}
  \item \ollabel{defn:term-var} $x$ ચલ હોય તો $x$ પદ છે.
  \item \ollabel{defn:term-abs} $x$ ચલ અને $M$ પદ હોય તો
    $(\lambd[x][M])$ પદ છે.
  \item \ollabel{defn:term-app} $M$ અને $N$ બંને પદો હોય તો
    $(MN)$ પદ છે.
\end{enumerate}
\end{defn}

કોઈ પદ~$(\lambd[x][M])$ \olref{defn:term-abs} મુજબ રચાયું હોય
તો તેને \emph{અમૂર્તીકરણ}નું પરિણામ કહીએ છીએ અને તેમાં
$\lambd[x]$માંના~$x$ને \emph{!!{parameter}} કહીએ છીએ. એ જ રીતે,
\olref{defn:term-app} મુજબ રચાયેલું પદ~$(MN)$
\emph{અનુપ્રયોગ}નું પરિણામ છે.

ઉપર વ્યાખ્યાયિત પદોમાં બધા જરૂરી કૌંસ છે. આ ક્યારેક ભારે પડે છે;
દાખલા તરીકે પદ
$(\lambd[x][((\lambd[x][x])(\lambd[x][(xx)]))])$ જુઓ. કૌંસ
છોડી શકાય તે માટે આપણે આગળ નિયમો આપીશું. જોકે, આ ઔપચારિક
વ્યાખ્યાથી \olref{defn:term} મુજબ પદ કઈ રીતે રચાયું છે તે નક્કી
કરવું સરળ બને છે. દાખલા તરીકે,
$(\lambd[x][((\lambd[x][x])(\lambd[x][(xx)]))])$ પદની રચનાનું
છેલ્લું પગલું અમૂર્તીકરણ જ હોવું જોઈએ, જેમાં !!{parameter}~$x$ છે.
તે પદ~$((\lambd[x][x])(\lambd[x][(xx)]))$માંથી અમૂર્તીકરણ વડે
બને છે; આ બીજું પદ બે પદોનો અનુપ્રયોગ છે. તે બંને પદો પોતપોતે
અમૂર્તીકરણથી બને છે, અને આ રીતે આગળ જઈ શકાય છે.

\begin{prob}
$(\lambd[g][((\lambd[x][(g (x x))])
  (\lambd[x][(g (x x))]))])$ની રચના વર્ણવો.
\sourcecorrection{OLLAM-010}{સ્થિર અંગ્રેજી મૂળના આ અભ્યાસપ્રશ્નમાં
બે અમૂર્તીકરણોનો અનુપ્રયોગ બહારના કૌંસ વિના લખાયો છે. અહીં
વ્યાખ્યાના~$(MN)$ નિયમ અનુસાર તેમની આસપાસ કૌંસ ઉમેર્યા છે.}
\end{prob}

\end{document}
